\documentclass[10pt,a4paper]{article}
\usepackage{fontspec}
\usepackage[english,russian]{babel}
\babelfont{rm}{PT Serif}
\babelfont{sf}{PT Sans}
\babelfont{tt}[Scale=0.88]{PT Mono}
\usepackage[a4paper,top=20mm,bottom=22mm,left=20mm,right=20mm]{geometry}
\usepackage{amsmath}
\usepackage{booktabs}
\usepackage{array}
\usepackage{enumitem}
\usepackage{titlesec}
\usepackage{caption}
\usepackage{tikz}
\usetikzlibrary{arrows.meta}
\usepackage{xcolor}
\usepackage[hidelinks]{hyperref}
\usepackage{microtype}

\setlength{\parindent}{0pt}
\setlength{\parskip}{3pt plus 1pt}
\setlength{\emergencystretch}{4em}

\titleformat{\section}{\normalfont\fontsize{11.5}{14}\bfseries}{\thesection.}{0.5em}{}
\titlespacing*{\section}{0pt}{8pt plus 2pt minus 2pt}{3pt}
\titleformat{\subsection}{\normalfont\fontsize{10}{12}\bfseries}{\thesubsection.}{0.5em}{}
\titlespacing*{\subsection}{0pt}{5pt plus 2pt minus 1pt}{2pt}

\setlist[enumerate]{leftmargin=2.0em,itemsep=0.5pt,topsep=2pt,parsep=0pt}
\setlist[itemize]{leftmargin=1.5em,itemsep=0.5pt,topsep=2pt,parsep=0pt}
\captionsetup{font=small,labelfont=bf,labelsep=period,skip=3pt}
\newcolumntype{L}[1]{>{\raggedright\arraybackslash}p{#1}}
\setlength{\textfloatsep}{7pt plus 2pt minus 2pt}
\setlength{\floatsep}{8pt plus 2pt minus 2pt}

\pagestyle{plain}

\tikzset{
  box/.style={draw,rectangle,thick,align=center,inner sep=2.5pt,font=\footnotesize,minimum width=27mm,minimum height=8mm},
  zone/.style={draw,dashed,rounded corners=2pt},
  zl/.style={font=\scriptsize\itshape,anchor=north west,inner sep=2pt},
  lbl/.style={font=\scriptsize,align=center,fill=white,inner sep=1pt},
  arr/.style={-{Latex[length=1.8mm]},thick},
}

\newcommand{\cat}{\,\Vert\,}

\begin{document}

\begin{center}
{\fontsize{16}{19.5}\selectfont\bfseries Swarma Chain: личности, имена и ценность\\[2pt] для роёв автономных агентов на одних хэш-функциях\par}
\vspace{5pt}
{\normalsize Ячейки личностей с предротацией, проверяемые аттестации,\\ собственный майнинг и штраф поздней ветки\par}
\vspace{8pt}
{\normalsize Алексей Григорьев, Иван Опарин\\
Relux Works\\
\href{https://github.com/relux-works}{github.com/relux-works}\par}
\vspace{3pt}
{\normalsize Октябрь 2026, версия 1.0\par}
\end{center}

\vspace{4pt}

\begin{list}{}{\leftmargin=0.6cm\rightmargin=0.6cm}
\item\small
\textbf{Аннотация.}
Агенты разных владельцев встречаются на границе проектов: ревьювер приходит в чужой рантайм, рабочая нода просится в открытый проект, работа оплачивается между владельцами.
Принимающей стороне нужно узнать участника, проверить его историю и рассчитаться с ним без доверия базе другого проекта и реестру другого владельца.
Swarma Chain даёт для этого общее состояние.
Личность в нём задаёт расходуемая ячейка с текущими ключами, хэшем следующих ключей и ссылкой на родителя, и состояние ячеек входит в консенсус наравне с деньгами и именами.
Аттестации об исходе работы подписываются пачками, а корень состояния личностей и имён лежит в каждом заголовке, поэтому субъект носит доказательства с собой, и его историю проверяют без доверия индексатору.
Консенсус опирается только на хэш-функции: обязательства SHA-256$\cat$SHAKE256, подписи парой SLH-DSA-128s из двух семейств хэшей, случайность из хэшей блоков и луковиц хэшей.
Кривых, пейрингов и решёток в нём нет, поэтому прорыв в их криптоанализе не подделывает ни личностей, ни платежей.
Доказательство работы на собственной конфигурации RandomX отвечает только за порядок записей.
Штраф поздней ветки требует от задержанной ветки кратного перевеса работы, который растёт с задержкой, и через 24 часа исключает реорганизации: при ожидании получателя в 3 часа двойная трата требует больше 87\,\% хэшрейта.
Мы задаём модель угроз, описываем устройство и разбираем каждую гарантию с допущениями и пределом.
Параметры оценены расчётом по явным кандидатным значениям.
\end{list}

\section{Введение}

Агентом мы называем процесс с языковой моделью, который вызывает инструменты, меняет файлы и отправляет сообщения.
Владелец отвечает за свои машины, своих агентов и оплату моделей.
Рантайм (Swarma Runtime) исполняет рой агентов на машинах одного владельца: у каждого агента своя учётка ОС, закрытые ключи держит служба ключник, а каждое действие проверяется по гранту \cite{swarmapaper}.
Протокол (Swarma Protocol) позволяет агентам независимых владельцев вести общий проект над единой средой исполнения, с журналом на кворуме и независимым ревью каждого кандидата.
Swarma Chain образует третий слой платформы: цепочку с общим состоянием для людей и агентов.

Внутри одной установки владелец знает, каких агентов он создал и какие роли им поручил.
На границе владельцев этого знания нет.
Внешний ревьювер приходит в чужой проект без локальной предыстории, и если каждый проект узнаёт его только по своему списку, работа в одном проекте ничего не говорит следующему.
Реестр владельца, который ведёт ключи его агентов, решает задачу внутри установки, но владелец администрирует его сам и может показать разным проверяющим разные истории одной личности.
Проектам нужно рассчитываться с участниками других владельцев, людям и агентам нужны имена, публичной точке присоединения нужна цена заявки против множественных личностей, а выбору ревьювера нужна случайность, которую никто не подбирает после появления кандидата.

Каждую из этих задач по отдельности решают известные средства: журнал прозрачности ключей \cite{ct2,coniks}, сервер имён, платёжная сеть \cite{bitcoin}, служба репутации.
Рою нужна их связка.
Имя указывает на личность, личность выдаёт и получает аттестации, аттестация действительна, если личность выдавшего была жива на высоте записи, а оплата записей и расчёты идут в одной валюте.
Если эти части держат разные операторы, проверка истории снова требует доверия каждому из них.
Swarma Chain держит их в одном состоянии, которое проверяет любой полный узел.

Второе требование касается срока жизни.
История участника должна сохранять доказательную силу годами, а системы ИИ ускоряют работу над математикой, в том числе криптоанализ.
Кривые и решётки могут оказаться слабее, чем считалось, а квантовый компьютер решает задачу дискретного логарифма, на которой держатся подписи на кривых \cite{shor}.
Хэш-конструкции остаются самой консервативной опорой.
Поэтому консенсус Swarma Chain опирается только на хэш-функции, и за это мы платим размером подписей и пропускной способностью.

Работа опирается на четыре понятия.
Личностью называется устойчивая криптографическая сущность человека, агента или организации: от её имени участник действует, и к ней относится история его работы.
Ячейкой личности называется расходуемая запись в состоянии цепочки с текущими ключами, хэшем следующих ключей и ссылкой на родителя (раздел~\ref{sec:identity}).
Аттестацией называется подписанное свидетельство об исходе работы участника, а выдавшим называется рантайм проекта, который подписал её своей личностью (раздел~\ref{sec:attest}).

Вклад работы:
\begin{enumerate}
\item Модель системы и модель угроз для общего слоя личностей роя агентов, в которой противник может арендовать большинство хэшрейта, сломать класс структурированной криптографии, украсть ключи и выпускать новые личности.
\item Личность как ячейка цепочки: расходуемый выход с предротацией, внешним родителем и терминальным отзывом, состояние которого проверяют полные узлы, и переход от журнала ключевых событий в реестре владельца без смены идентификатора.
\item Криптографический профиль консенсуса на одних хэш-функциях: двойные обязательства, пара подписей SLH-DSA из двух семейств хэшей, пачки с одной подписью на корень и случайность без VRF и пороговых подписей.
\item Финальность без второго консенсуса: собственный майнинг RandomX без объединённого, штраф поздней ветки с отсечкой реорганизаций через 24 часа, checkpoint для новых узлов и майнинг как побочный продукт работы роя, где PoW-билеты допуска служат шарами.
\item Аттестации с доказательством у субъекта и репутация как неподвижная точка графа подписанных утверждений, которую проверяют без индексатора.
\item Анализ гарантий с допущениями и пределами и расчёт кандидатных параметров: пропускной способности, цены записи, роста цепочки заголовков, графика эмиссии и цены двойной траты.
\end{enumerate}

Подробное устройство платформы описывает книга о Swarma \cite{swarmabook}.
Если работа расходится со спецификацией модуля, права спецификация.

\section{Модель системы и модель угроз}
\label{sec:model}

\textbf{Участники.}
Таблица~\ref{tab:roles} перечисляет участников цепочки и то, что проверяющему приходится о каждом из них принять.
Агенты, люди и организации владеют личностями.
Рантайм проекта выдаёт аттестации от своей личности, и выдавать их вправе любой рантайм, у которого есть личность в цепочке.
Майнеры упорядочивают записи, полные узлы проверяют блоки и состояние, лёгкие клиенты проверяют заголовки и доказательства.
Индексатор вычисляет производное состояние: профили, деревья ответов, разрешение указателей и репутацию.

\begin{table}[t]
\centering\small
\caption{Участники Swarma Chain.}
\label{tab:roles}
\footnotesize
\begin{tabular}{L{3.2cm}L{6.5cm}L{5.9cm}}
\toprule
Участник & Что делает & Что проверяющий принимает о нём \\
\midrule
Агент, человек, организация & Владеет личностью; подписывает события личности и пачки, агент через ключник & Ничего сверх подписи и состояния ячейки \\
Владелец & Служит родителем ячеек своих агентов; хранит ключи своей ячейки и внешний ключ восстановления & Сохранность ключей владельца \\
Рантайм проекта (выдавший) & Выдаёт аттестации пачками от своей личности & Вес по истории; истинность свидетельства не принимается \\
Майнер & Упорядочивает записи доказательством работы & Порядок свежих записей; содержание проверяется отдельно \\
Полный узел & Проверяет блоки, траты, состояние личностей и имён, конверты записей & Подлинность checkpoint при старте \\
Лёгкий клиент & Проверяет заголовки от checkpoint, фильтры и доказательства & Подписантов checkpoint \\
Индексатор & Вычисляет производное состояние и репутацию по опубликованным правилам & Ничего: вывод пересчитывается \\
Архивист, сервисный узел & Хранит и раздаёт старые сегменты и медиа & Наличие копий, без гарантии \\
\bottomrule
\end{tabular}
\end{table}

\textbf{Допущения.}
\begin{enumerate}[label=Д\arabic*.]
\item Хэш-функции: хотя бы одно из семейств SHA-2 \cite{fips180} и SHA-3 \cite{fips202} сохраняет стойкость к коллизиям и вторым прообразам на используемых длинах, и стойкость каждого варианта SLH-DSA \cite{fips205} сводится к свойствам его хэш-функции.
\item Хэшрейт: противник может арендовать большинство хэшрейта на ограниченное время; гарантии финальности выражены через его долю $\alpha$ и время ожидания получателя (раздел~\ref{sec:finality}).
\item Сеть: честный узел связан хотя бы с одним честным пиром и получает честные блоки с задержкой, малой по сравнению с 10 минутами.
\item Время: у узла есть монотонные часы, а системное время нужно ему только для проверки меток блоков с допуском 300 секунд в будущее.
\item Старт: новый узел получает подлинный checkpoint, который подписывают хэш-ключи сопровождающих и который публикуется в релизе и на зеркалах.
\item Ключи владельца: ключи ячейки владельца хранятся отдельно от мастер-секрета установки, а внешний ключ восстановления не живёт на повседневных устройствах; противник не получает их оба, и владелец не теряет их оба.
\end{enumerate}

\textbf{Противники.}
\begin{enumerate}[label=П\arabic*.]
\item Криптоаналитик: квантовый компьютер или прорыв в математике ломает кривые, пейринги или решётки, а в худшем случае одно из двух семейств хэшей.
\item Арендатор большинства: противник временно держит больше половины хэшрейта и пытается переупорядочить свежие записи, задержать запись или дважды потратить недавний платёж.
\item Вор ключей: противник получил текущие ключи ячейки или мастер-секрет ключника установки.
\item Двуличный владелец: владелец показывает разным проверяющим разные истории одной личности.
\item Фабрика личностей: противник создаёт много личностей и аттестует ими себя \cite{sybil}, или несколько выдавших сговариваются.
\item Нечестный индексатор: сервер искажает производное состояние, скрывает записи или подставляет другую формулу.
\item Сетевой противник: изолирует узел eclipse-атакой \cite{eclipse} или задерживает распространение блоков.
\item Спамер: заливает блоки корректно подписанными и оплаченными записями.
\end{enumerate}

\textbf{Цели.}
\begin{enumerate}[label=Ц\arabic*.]
\item Состояние денег, имён и личностей меняется только подписью, которую проверяют по самому состоянию.
\item Кража текущих ключей не перехватывает личность, а кража мастер-секрета установки не лишает владельца личностей его агентов.
\item После созревания у всех честных узлов одна версия состояния каждой личности и каждого имени.
\item Проверяющий устанавливает подпись, включение и статус выдавшего любой аттестации без доверия индексатору и выдавшему.
\item Цена двойной траты растёт с ожиданием получателя, а через 24 часа реорганизация невозможна.
\item Каждая запись платит с первого байта, а платежи, имена и события личностей не вытесняются потоком социальных записей.
\item История, закреплённая в цепочке, сохраняет доказательную силу на дату наблюдения после поломки класса ключей, которыми подписаны сами записи.
\end{enumerate}

\textbf{Вне модели.}
Цепочка не устанавливает истинности аттестации: подпись относит свидетельство к выдавшему и правоты ревьювера не доказывает.
Доказательство контроля над доменом не доказывает честности владельца.
Хранение контента вне цепочки ничем не гарантировано: указатель в блоке файла не сохраняет.
Шифрование с открытым ключом в консенсус не входит, поэтому конфиденциальность содержимого задают сами участники, а открытые метаданные раскрывают рабочие связи.
Устойчивое большинство может часами задерживать записи, и от цензуры такого противника цепочка не защищает.

\section{Устройство: подслои и граница консенсуса}
\label{sec:overview}

Цепочка заверяет записи, как нотариус.
Всё ценное в ней самоаутентично: подписи, предротация личностей, обязательства.
Доказательство работы отвечает только за порядок записей.
Противник с большинством хэшрейта ничего не подделает, но может переупорядочить свежие записи, задержать запись или дважды потратить недавний платёж, и защита строится против этих трёх действий.

\textbf{Три подслоя.}
Swarma Chain разделяет цепочку, p2p-память и социальный граф.
Цепочка задаёт общий порядок записей и хранит деньги, имена, личности, социальные записи и указатели на контент.
P2P-память переносит тяжёлый контент, её основной транспорт BitTorrent, и гарантии хранения у неё нет.
Социальный граф из подписок, отметок и категоризаций направляет, что участники скачивают, раздают и архивируют.
Поэтому полный узел не хранит медиа сети, а указатель в блоке доказывает публикацию, но не доступность файла.

\textbf{Где проходит консенсус.}
Консенсус ведёт три вида состояния.
Для денег он согласует расходование UTXO.
Для имён он обеспечивает уникальность, владение, продление и передачу.
Для личностей он держит ячейки.
Если бы две несовместимые траты, два владельца одного имени или две версии одной личности могли одновременно считаться правильными, клиент не исправил бы расхождение собственным выбором.
Поэтому это состояние проверяют узлы, а корень хэш-дерева состояния личностей и имён лежит в каждом заголовке.

У остальных записей узел проверяет конверт: тип, размер, подпись пачки, оплату и статус выдавшего.
Смысл записей детерминированно выводят индексаторы по опубликованным правилам, а что показать читателю, решает клиент.
Генезис сразу поддерживает всё, что нужно социальному взаимодействию: сегмент данных, канонический типизированный конверт записи, пачки с одной подписью на корень, тариф по типу, обрезку старых сегментов и фильтры для лёгкого клиента.
Аттестация служит одним из типов записей того же механизма.
Неизвестный узлу тип принимается как данные по общему тарифу, а новый смысл задаёт соглашение индексаторов, поэтому социальный продукт растёт поверх консенсуса без форков.
Новых правил расходования денег такое расширение не даёт.
Социальные записи цепочка хранит как нетратимые типизированные выходы, и они не засоряют UTXO-набор; имена и ячейки личностей хранятся как расходуемые выходы.

\textbf{Блок.}
Блок состоит из тела и сегмента данных (рис.~\ref{fig:block}).
Тело несёт платежи, операции с именами и события личностей.
Сегмент данных несёт социальные записи и пачки аттестаций.
Заголовок фиксирует корни обеих частей, корень состояния личностей и имён и обязательство фильтра для лёгких клиентов.
Полный узел проверяет сегмент при приёме блока и вправе потом его выбросить: обрезка удаляет социальную историю и не трогает состояния денег, имён и личностей.
Интервал блока равен 60 секундам.
Точный формат записей задаёт спецификация цепочки.

\begin{figure}[t]
\centering
\begin{tikzpicture}
  \node[box,minimum width=140mm] (hdr) at (0,2.1) {заголовок: корень тела, корень сегмента данных, корень состояния личностей и имён, обязательство фильтра};
  \node[box,minimum width=62mm] (body) at (-4.3,0.6) {тело\\\scriptsize платежи, операции с именами, события личностей};
  \node[box,minimum width=62mm] (seg) at (4.3,0.6) {сегмент данных\\\scriptsize социальные записи, пачки аттестаций; узел вправе выбросить};
  \node[box,minimum width=62mm] (state) at (-4.3,-0.95) {состояние\\\scriptsize UTXO, имена, ячейки личностей};
  \node[box,minimum width=62mm] (batch) at (4.3,-0.95) {пачка\\\scriptsize дерево Меркла записей, парная подпись корня};
  \draw[arr] (hdr.south -| body.north) -- node[lbl,right=1pt]{корень} (body.north);
  \draw[arr] (hdr.south -| seg.north) -- node[lbl,right=1pt]{корень} (seg.north);
  \draw[arr] (body) -- node[lbl,right=1pt]{меняет} (state);
  \draw[arr] (batch) -- node[lbl,right=1pt]{лежит в} (seg);
  \draw[arr] (hdr.west) -| ++(-0.35,-1) |- node[lbl,pos=0.25,rotate=90]{корень состояния} (state.west);
\end{tikzpicture}
\caption{Блок Swarma Chain. Тело меняет состояние денег, имён и личностей, корень этого состояния лежит в заголовке. Сегмент данных несёт пачки, и узел вправе выбросить его после проверки; субъект аттестации хранит свой путь до корня пачки сам.}
\label{fig:block}
\end{figure}

\section{Криптографический профиль: только хэш-функции}
\label{sec:crypto}

Консенсус Swarma Chain опирается на один класс предположений: стойкость хэш-функций.
Для подписей и обязательств этого достаточно.
Шифрование с открытым ключом без структурированной математики не построить, но в консенсус оно не входит.
Таблица~\ref{tab:crypto} сводит выбор.

\begin{table}[t]
\centering\small
\caption{Криптографический профиль консенсуса.}
\label{tab:crypto}
\footnotesize
\begin{tabular}{L{3.5cm}L{4.2cm}L{4.1cm}L{3.4cm}}
\toprule
Что & Выбор & Зачем & Цена \\
\midrule
Обязательства: идентификаторы блоков, транзакций и личностей, деревья Меркла & SHA-256$\cat$SHAKE256, 64 байта, с разделением доменов & Коллизия требует сломать оба семейства, SHA-2 и SHA-3 & Идентификаторы вдвое длиннее \\
Подписи: деньги, имена, события личности, транзакции пачек & Пара SLH-DSA-SHA2-128s и SLH-DSA-SHAKE-128s, обязательны обе & Только хэш-предположения; поломка одного семейства не даёт подделки & 15\,712 байт подписи; подпись от 0,1 до 1~с, проверка за миллисекунды \\
Аттестации и социальные записи & Пачка: дерево Меркла, одна парная подпись на корень & Подпись делится на сотни записей & Путь $64\lceil\log_2 N\rceil$ байт \\
Ротация личности & Предротация: событие несёт хэш следующей пары ключей & Украденный или взломанный текущий ключ не перехватывает личность & 64 байта на событие \\
Адреса монет & Хэш условия траты, одноразовые адреса & Ключ скрыт до траты & \\
Случайность для выбора & Хэш блока и луковица хэшей & VRF и пороговые BLS держатся на кривых и пейрингах & \\
Доказательство работы & Собственная конфигурация RandomX & Нагрузка на CPU; ей угрожает железо, прорыв в математике её не ломает & \\
\bottomrule
\end{tabular}
\end{table}

\textbf{Обязательства.}
Каждое обязательство состоит из двух половин: SHA-256 и SHAKE256 с выходом 256 бит, вместе 64 байта, с разделением доменов.
Так устроены идентификаторы блоков, транзакций и личностей и узлы деревьев Меркла.
Коллизия обязательства требует одновременной коллизии в SHA-2 и в SHA-3, поэтому прорыв в одном семействе обязательства не подделывает.
Стойкость конкатенации к коллизиям близка к стойкости более сильной из двух функций и их сумме не равна: для итеративных конструкций многократные коллизии строятся почти по цене одной \cite{joux}.
Поэтому выигрыш профиля даёт независимость двух семейств, и длина выхода здесь роли не играет.

\textbf{Подписи.}
Каждая подпись консенсуса состоит из двух подписей SLH-DSA-128s \cite{fips205}: варианта SHA2 и варианта SHAKE.
Подпись действительна, только если верны обе.
SLH-DSA продолжает линию SPHINCS и SPHINCS$^+$ \cite{sphincs,sphincsplus}: схема не хранит состояния, а её стойкость сводится к свойствам хэш-функции.
Набор 128s даёт подпись в 7856 байт при открытом ключе в 32 байта и соответствует категории стойкости 1 по классификации NIST, так что пара весит $2\times7856=15\,712$ байт подписи и 64 байта ключа.
Подпись занимает от 0,1 до 1 секунды, проверка занимает миллисекунды.

Три варианта мы сознательно не берём.
Пара SLH-DSA-256s весит около 60~КБ, а от прорыва в математике защищают два независимых семейства, и длина параметров тут не помогает.
У ключей WOTS, XMSS и LMS \cite{xmss,rfc8391,rfc8554} есть состояние, и оно ломается при восстановлении из резервной копии и при повторной подписи; для таких схем NIST задаёт отдельные требования к управлению состоянием \cite{sp800208}.
384-битные хэши не нужны, потому что слабое место не в размере выхода: при росте опасений сначала добавляют раунды.

\textbf{Адреса и случайность.}
Адрес монет равен хэшу условия траты, и адреса одноразовые, поэтому открытый ключ скрыт до траты.
Случайность для выбора дают хэш блока и луковица хэшей.
Луковицей хэшей называется цепочка хэшей, вершину которой участник публикует заранее, а в каждом раунде раскрывает следующий прообраз.
VRF \cite{vrf} и маяки на пороговых подписях BLS держатся на кривых и пейрингах и в профиль не входят.

\textbf{Ключи агентов.}
Ключник рантайма выводит ключи ячейки агента из мастер-секрета установки хэш-функцией SHAKE256 от пути личности и номера ротации.
Ключи для каналов рантайма выводятся отдельно, по SLIP-0010 \cite{slip10} для Ed25519.
Закрытый ключ личности агенту не выдаётся: ключник подписывает по его запросу.

\textbf{Корни доверия вне консенсуса.}
Остальная Swarma стоит на структурированной криптографии: SSH-сертификаты Ed25519, подписи журнала протокола и квитанций на Ed25519, ключ оператора на телефоне (P-256 в защищённом чипе), ключи FIDO2, TLS и подписи git.
Если кривые сломаются, противник подделает подписи этими ключами.
Дизайн признаёт этот риск и ограничивает его четырьмя шагами.
Первый шаг: головы журналов протокола и реестров владельцев закрепляются в цепочке, и запись, наблюдённая до поломки, сохраняет доказательную силу на дату наблюдения.
Второй шаг: долгоживущие корни, начиная с ячейки владельца, подписывают хэш-подписями, а ключи на кривых остаются короткоживущими листьями под ними, так что противник подделает лишь живые листья на время их действия и не выпустит новых.
Третий шаг: сеансовый слой следует за вендорами и использует гибридный обмен ключами с ML-KEM \cite{fips203}.
Четвёртый шаг: политика отвергает поломанный класс ключей, а личности ротируются через предротацию, потому что их следующие ключи спрятаны за хэшем.

\section{Личность как ячейка}
\label{sec:identity}

\textbf{Ячейка.}
Ячейка личности хранит три поля.
Поле текущих ключей хранит пару открытых ключей SLH-DSA-128s в вариантах SHA2 и SHAKE.
Поле следующих ключей хранит 64-байтное обязательство на пару, которая заменит текущую.
Поле родителя указывает на ячейку владельца, а у самого владельца на внешний ключ восстановления.
Идентификатор личности равен хэшу первой записи ячейки, и ротации его не меняют.

\textbf{Три пути траты.}
Ячейку меняют только траты (рис.~\ref{fig:cell}).
Ротация тратит ячейку следующими ключами: их хэш должен совпасть с записанным обязательством, обе подписи пары обязательны, а новая ячейка несёт новые текущие ключи и обязательство на очередные.
Трата родителем подписана текущими ключами родительской ячейки и задаёт новые текущие ключи и новое обязательство на следующие.
Отзыв тратит ячейку в терминальную ячейку без преемника, и личность больше не подписывает.
Собственная ротация не может изменить ссылку на родителя.
Любая трата действует с момента публикации: подделать её нельзя, поэтому задерживать её цензурой бесполезно.

Текущие ключи ячейку не тратят.
Ими подписываются пачки аттестаций и социальные записи.
Так ключ повседневной подписи отделён от права менять личность.
Предротация, как в KERI \cite{keri}, прячет следующие ключи за хэшем, и их не вычислить ни из открытых текущих ключей, ни из подписей текущими.
Даже если класс текущих ключей окажется слабым, ротацию через заранее объявленные следующие ключи противник не перехватит.

\begin{figure}[t]
\centering
\begin{tikzpicture}
  \node[box,minimum width=46mm] (c0) at (0,0) {ячейка, версия $v$\\\scriptsize текущие $K_v$; следующие $H(K_{v+1})$; родитель $p$};
  \node[box,minimum width=50mm] (c1) at (9.6,1.45) {версия $v{+}1$: ротация\\\scriptsize текущие $K_{v+1}$; следующие $H(K_{v+2})$; родитель $p$};
  \node[box,minimum width=50mm] (cp) at (9.6,0) {версия $v{+}1$: трата родителем\\\scriptsize текущие $K^{*}$; следующие $H(K^{*}_{+})$; родитель $p$};
  \node[box,minimum width=50mm] (ct) at (9.6,-1.45) {терминальная ячейка\\\scriptsize преемника нет};
  \node[box,minimum width=46mm] (rec) at (0,-1.45) {пачки и социальные записи\\\scriptsize подписаны текущими $K_v$};
  \draw[arr] (c0.east) -- node[lbl,sloped,above=1pt,pos=0.55]{подпись $K_{v+1}$, хэш совпал} (c1.west);
  \draw[arr] (c0.east) -- node[lbl,above=1pt,pos=0.5]{подпись ключами ячейки $p$} (cp.west);
  \draw[arr] (c0.east) -- node[lbl,sloped,below=1pt,pos=0.55]{отзыв} (ct.west);
  \draw[arr,dashed] (c0) -- node[lbl,right=1pt]{не тратят} (rec);
\end{tikzpicture}
\caption{Три пути траты ячейки личности. $H$ обозначает 64-байтное обязательство. Собственная ротация сохраняет ссылку на родителя $p$, поэтому путь траты родителем открыт в каждой версии.}
\label{fig:cell}
\end{figure}

\textbf{Владелец и восстановление.}
Ячейка владельца служит родителем ячеек его агентов.
Её ключи не выводятся из мастер-секрета установки и хранятся отдельно, на телефоне под защитой его чипа или на внешнем носителе.
Плановая смена мастер-секрета идёт через новый корень: агенты регистрируются под ним заново, их ячейки ротируются на новые ключи, а непрерывность держит ячейка владельца.
При компрометации мастер-секрета предротация не помогает: следующие ключи агентов выводятся из того же секрета, и вор знает их тоже.
Восстановление идёт вторым путём траты.
Ячейка владельца тратит ячейку каждого агента и задаёт ключи нового корня и обязательство на следующие.
Вор этих ключей не знает, поэтому после такой траты ячейка ему недоступна, даже если он успел провернуть её сам.
Записи, которые вор подписал в промежутке, владелец оспаривает подписанными возражениями (раздел~\ref{sec:attest}).
Если потеряны и ключ владельца, и ключ восстановления, личность не восстанавливается; этого система не обещает.

\textbf{Отзыв и записи.}
Записи, подписанные ключами ячейки после высоты её отзыва или ротации, недействительны, и это следует из корня состояния на высоте записи.
Записи, сделанные раньше, остаются действительными, и пересчёта прошлой истории нет.
Вывод учётки агента из работы и отзыв личности различаются: постоянного агента можно отключить в рантайме и позднее вернуть с той же личностью, а отзыв выполняется отдельно.
Иначе обслуживание машины могло бы уничтожить историю.

У каждого агента своя ячейка, поэтому его отзывают, не трогая остальных.
Ссылка на родителя объединяет происхождение личностей одного владельца и сохраняет отдельные ключи.
Личности под одной ячейкой владельца относятся к одному административному домену: разные ключи двух агентов не показывают, что за ними стоят разные владельцы.

\textbf{До цепочки: журнал ключевых событий.}
Рантайму личности нужны независимо от того, доступна ли цепочка, поэтому идентификатор самосертифицирующий \cite{keri}: он равен хэшу первой записи журнала ключевых событий.
Первая запись объявляет текущие ключи, обязательство на следующие и родителя.
Ротации и отзывы записываются следующими записями, и каждая подписана по правилам ячейки.
Журнал ведёт реестр владельца: он служит единственным авторитетным писателем, хранит журнал в git для аудита и выдаёт подписанные снимки не старше 60 секунд.
Проверить историю от первой записи может любой, кому журнал показан.
При публикации журнала в цепочку ячейка получает тот же идентификатор и последнее состояние, и ссылки на личность, накопленные до этого, остаются действительными.

Реестр одного владельца не защищает от раздвоения журнала.
Свидетель на узле доверия ловит откат и раздвоение в пределах сохранённых контрольных точек, но узлом администрирует сам владелец, и показать разным проверяющим разные истории он может.
Журналы прозрачности \cite{ct2,coniks,sigstore} делают такое раздвоение обнаружимым при наличии аудиторов.
В цепочке состояние ячейки одно у всех честных узлов, а после 24 часов его не меняет и реорганизация.

\textbf{Доказательство идентичности.}
Доказательство идентичности связывает личность с внешней сущностью в обе стороны.
Владелец подписывает заявление ключами своей ячейки и публикует у внешнего провайдера артефакт с идентификатором личности: TXT-запись домена под DNSSEC, файл well-known по HTTPS или ключи аккаунта на хостинге кода; публикации в социальных платформах образуют необязательный класс.
Внешний запрос недетерминирован, поэтому подтверждение требует, чтобы пороговое число валидаторов записало одинаковый хэш свидетельства в одном окне наблюдения.
Расхождение означает, что подтверждения нет.
Подтверждённое наблюдение публикуется в цепочке как аттестация о личности, с хэшем свидетельства и датой, и поздняя недоступность провайдера его не отменяет.
Классы доказательства весят по-разному: домен под DNSSEC весит больше публикации в социальной платформе.
Ни один класс не свидетельствует о честности владельца.

\section{Имена}
\label{sec:names}

Имя хранится как расходуемый выход с истечением и продлением и указывает на идентификатор личности, поэтому смена ключей не меняет того, на что оно указывает.
Регистрация идёт в два шага, обязательство и раскрытие, и имя вступает в силу через 24 часа.
Обязательство скрывает имя до раскрытия, поэтому тот, кто увидел заявку, не опередит её своей.
Срок созревания совпадает с отсечкой реорганизаций (раздел~\ref{sec:finality}): после него перехватить имя реорганизацией невозможно.

Цена регистрации зависит от длины имени и следует за спросом.
Плата делится пополам: одна часть сжигается, другая распределяется майнерам по блокам следующих суток.
Так плата за имя поддерживает безопасность и не собирает комиссии в отдельных особенно выгодных блоках.
Изучение Namecoin показало, что уникальности имён мало без сдерживания сквоттинга \cite{namecoinstudy}; цена по длине и спросу и истечение без продления служат этому сдерживанию.
ENS \cite{ens} решает задачу имён контрактом поверх Ethereum, а в Swarma Chain имя входит в состояние консенсуса и в тот же корень, что и личности.

\section{Аттестации и репутация}
\label{sec:attest}

\textbf{Состав.}
Аттестацию выдаёт рантайм проекта после ревью или приёмки работы.
Агент выступает субъектом свидетельства, проект выступает выдавшим.
Минимальный состав аттестации канонический: личность агента, роль (продьюсер, ревьювер или верификатор), ссылка на проект, хэш кандидата или ревизии, версии контракта задачи и политики, исход (\emph{принято}, \emph{отклонено} или \emph{не установлено}), высота и время, личность выдавшего и подпись пачки.
Хэш различает рассмотренные объекты, а версии контракта и политики сохраняют условия оценки.
Поэтому запись о прежнем исходе не становится свидетельством об изменённом коде или о новых условиях.

\textbf{Пачки.}
Выдавший собирает аттестации в пачку, строит из них дерево Меркла \cite{merkle89} и подписывает корень одной парной подписью в транзакции пачки.
Подпись в 15\,712 байт делится на сотни записей, а каждая запись добавляет к цене пачки свои байты и путь в дереве.
Цена аттестации равна тарифу записи, то есть базовой плате за байт, которая сжигается.
Отдельной платы за свидетельство нет.
От спама защищает вес выдавшего: свидетельства выдавших без веса почти ничего не дают.

\textbf{Проверка без чужой базы.}
Субъект носит доказательства с собой: свидетельство, путь Меркла, транзакцию пачки и заголовок блока.
Проверяющий выполняет четыре шага.
Подпись: пачку подписали ключи, текущие у выдавшего на высоте записи.
Включение: путь от свидетельства сходится к корню пачки, а пачка лежит в блоке под заголовком.
Статус: ячейка выдавшего на этой высоте жива, что видно по корню состояния личностей в заголовке.
Вес: он считается по опубликованной формуле.
Первые три шага не требуют доверия ни индексатору, ни выдавшему, поэтому репутацию субъекта можно проверить, даже если индексаторы исчезли.

\textbf{Открытые записи и обязательства.}
Аттестации по умолчанию открыты.
В цепочку попадают только хэши и метаданные, а код, промпты, данные проекта и секреты в неё не попадают.
Открытая запись всё же раскрывает, какой агент с каким проектом работал, поэтому приватный проект публикует только обязательство, то есть хэш аттестации.
Полное свидетельство получает субъект и раскрывает его тому, кому хочет, а проверяющий сверяет его с обязательством.
Публичный индексатор учитывает только открытые записи.
Обязательство скрывает содержимое, но не факт записи и не выдавшего.

\textbf{Репутация.}
Репутацией называется функция над аттестациями, которую детерминированно вычисляют индексаторы.
Индексатор считает её как неподвижную точку графа подписанных утверждений, по схеме EigenTrust \cite{eigentrust}.
Каждая личность в этом графе и выдаёт, и получает, потому что агенты выдавшего сами работают в чужих проектах.
Вес втекает в граф от затравки, а затравку задают классы доказательства личности.
Кластер личностей, которые аттестуют друг друга, получает ровно тот вес, который втекает в него снаружи.
Вес выдавшего зависит от класса доказательства его личности, его возраста и того, подтверждали ли его вердикты независимые выдавшие; вердикты, которые независимые выдавшие опровергли, вес снижают.
Оси продьюсера, ревьювера и верификатора считаются раздельно, а старые свидетельства весят меньше новых.
Исход \emph{не установлено} нейтрален и репутации не прибавляет и не отнимает.

Формула и её константы публикуются вместе с индексатором, и конкретные значения затухания и начальных весов классов задаёт эта публикация.
Проверяющий может выбрать собственную затравку и вычислить репутацию от своих доверенных точек; скрытого списка доверенных выдавших нет.
Детерминированность означает, что при одинаковых записях и одинаковой формуле получается одинаковый результат, и любой может пересчитать индексатор.
Правильности самой формулы она не доказывает.

\textbf{Споры.}
Спором называется подписанное возражение без арбитража.
Субъект публикует возражение со ссылкой на аттестацию, выдавший может опубликовать поправку или отзыв своей записи, а индексаторы показывают свидетельство и возражения рядом.
Оспоренное \emph{отклонено} весит меньше, и полный вес ему возвращает только повторное подтверждение выдавшего.
Арбитра в цепочке нет: каждая сторона отвечает своей подписью и своей историей.

\textbf{Деньги и репутация.}
Дизайн не связывает деньги и репутацию.
Аттестация стоит только тариф записи, оплату работы задаёт политика проекта, а цепочка даёт валюту расчётов.
Репутацию нельзя купить, заложить или потерять как залог, и монета не начисляется за число аттестаций.
Связь между ними допустима только там, где понятны её цена и последствия, и это открытый вопрос дизайна.

\textbf{Допуск остаётся у проекта.}
Переносится история подписанных утверждений, а право работать выдаёт принимающий проект по своей политике.
Например, работе нового агента проект назначает двух-трёх независимых ревьюверов, а работе агента с историей одного; эти числа иллюстрируют политику и порогов не задают.
Репутация не создаёт отсутствующей независимости: личности одного владельца остаются одним административным доменом, а общее семейство модели остаётся источником общих слепых зон.

\section{Ценность: деньги, плата и эмиссия}
\label{sec:value}

\textbf{Деньги.}
Монеты лежат в UTXO, адрес равен хэшу условия траты, и адреса одноразовые.
Swarma Chain служит слоем расчётов и личностей: подписанных объектов в блоке мало (раздел~\ref{sec:params}), поэтому расчёты медленные, и для мелких покупок цепочка не предназначена.
Чаевые автору служат нативной записью.
Клиент или спонсор может оплачивать записи пользователей, и спонсор виден в записи.

\textbf{Размер блока и плата.}
Правило размера блока целиком зашито в генезис (табл.~\ref{tab:blocksize}), без голосований и без хардфорков ради размера.
Правило учитывает два урока Monero: короткая медиана позволяла экспоненциальный рост размера, а место до зоны без штрафа было почти бесплатным, и его дёшево заливали.
Базовая плата за байт меняется на $\pm12{,}5\,\%$ за блок по отклонению от цели, как в EIP-1559 \cite{eip1559}, и сжигается.
Минимальная плата привязана к награде, поэтому заливка стоит денег с первого байта.
Рост цели требует устойчиво высокой платы, так что короткий всплеск спроса цель не двигает.
У сегмента данных свой предел, поэтому платежи, имена и события личностей не вытесняются потоком записей.
Срока обслуживания правило не обещает: каждая транзакция ждёт места в блоке.

\begin{table}[t]
\centering\small
\caption{Правило размера блока, зашитое в генезис. Числа служат кандидатами для моделирования.}
\label{tab:blocksize}
\footnotesize
\begin{tabular}{L{2.4cm}L{13.6cm}}
\toprule
Часть & Значение \\
\midrule
Старт & Цель $T_0=512$~КБ на 60-секундный блок; блок не больше $2T$ \\
Плата & Базовая плата за байт меняется на $\pm12{,}5\,\%$ за блок по отклонению от цели и сжигается; полный по цели блок стоит не меньше $R/16$, где $R$ награда за блок \\
Рост & Цель растёт, только когда базовая плата выше двойного минимума, и не быстрее $\times2$ в год \\
Спад & Цель падает, только когда блоки полупустые при минимальной плате, и не быстрее $\times\frac12$ в год; ниже $T_0$ не опускается \\
Потолок & Цель не выше $4~\text{МБ}\times2^{y/4}$, где $y$ возраст сети в годах \\
Состав & Вес включает все байты, подписи без скидки; социальный сегмент не больше половины веса; пересчёт на каждом блоке в целых числах \\
\bottomrule
\end{tabular}
\end{table}

\textbf{Эмиссия.}
Комиссии почти бесплатного текста дают слабую основу бюджета безопасности.
Если поднять их, публикации подорожают, а если держать низкими без другого дохода, вычислительная защита недополучит оплату.
Поэтому выбрана хвостовая эмиссия по образцу Monero \cite{tail}: награда плавно уменьшается вдвое каждые три года от 64 условных монет за блок и с пятнадцатого года остаётся постоянной, 2 монеты за блок.
Числа графика служат кандидатами, а расчёт приведён в разделе~\ref{sec:params}.
Эмиссия оплачивает продолжение работы сети и доходности держателю не обещает.
Запуск честный: без премайна и без эмитента, доля фонда разработки равна нулю.
Майнер получает награду, приоритетные чаевые и половину платы за имена; базовая плата за байт сжигается.

\textbf{Сервисные узлы.}
Сервисным узлом называется оператор полезной сетевой функции: индексирования, раздачи медиа, разрешения имён или хранения старых сегментов.
Власти над блоками у него нет, и консенсус не вводит ни мастернод, ни залоговых барьеров, ни кворумов.
Оплату сервисных узлов дизайн связывает с проверяемой работой.
Случайный пир запрашивает кусок сегмента, разрешение имени или подписанный хэш состояния индекса и проверяет ответ, а узел, не прошедший проверку, оплаты не получает.
Параметры этого механизма составляют открытое пространство дизайна: доля награды, наличие залога, частота проверок и ответственность за систематическую ложь.

\section{Консенсус и майнинг}
\label{sec:consensus}

\textbf{Собственный RandomX.}
Майнеры работают на собственной конфигурации RandomX \cite{randomx}, которая нагружает CPU, со своей солью, своей деривацией ключа по высоте блока и своими размерами программ.
Угрозу для такой функции создаёт железо, а прорыв в математике её не ломает.
Своя конфигурация затрудняет немедленный перенос чужих настроек майнинга, но ботнетов и арендованных CPU не устраняет.
Поэтому безопасность не оценивают по одному отсутствию известных ASIC: доли пулов наблюдаются, а параметры в первые два-три года сети меняются на плановых хардфорках.

\textbf{Без объединённого майнинга.}
Объединённый майнинг \cite{mergedmining} исключён по трём причинам.
Крупный пул сети-донора атаковал бы молодую цепочку бесплатно, потому что он и так майнит.
Monero сам держится на кривых, поэтому при их поломке его ценность и хэшрейт упадут ровно тогда, когда хэш-профиль Swarma Chain должен окупиться, и заёмная безопасность исчезнет в самый нужный момент.
Без второй дорожки консенсус проще, а раздача монет целиком своя.

\textbf{Сложность.}
Небольшая PoW-сеть чувствительна к переключающимся майнерам: они приходят на низкую сложность и уходят на высокую, а длинное окно пересчёта растягивает период после ухода хэшрейта.
Сложность пересчитывает ASERT \cite{asert} в целочисленной форме с якорем в генезисе, без скользящего окна.
В идеализированной записи цель блока высоты $h$ с меткой времени $t_h$ равна
\[
T_h = T_a\cdot 2^{\,(t_h - t_a - 60\,\text{с}\cdot(h-h_a))/\tau},
\]
где $(h_a,t_a,T_a)$ задаёт якорь, а $\tau$ называется полупериодом: каждые $\tau$ отставания от расписания цель удваивается, то есть сложность падает вдвое.
Кандидат полупериода равен двум часам.
Метки времени строго монотонны, допуск в будущее равен 300 секундам, минимальной сложности на основной сети нет.

\textbf{Рой и есть майнеры.}
Хосты Swarma Runtime по желанию майнят простаивающим CPU: под отдельной системной учёткой, с низким приоритетом и лимитом cgroup, с сетью только для p2p цепочки.
Тысяча таких хостов даёт базу хэшрейта, которая не уходит при росте сложности.
Майнеры объединяются в децентрализованный пул в духе P2Pool \cite{p2pool} с порогом выплат, потому что мелкие выходы съела бы цена их траты.
Шаром называется решение облегчённой цели, которым майнер доказывает пулу свою работу.

\textbf{PoW-билеты допуска.}
Публичная точка присоединения протокола принимает заявки от незнакомых нод, поэтому заявка стоит работы \cite{dworknaor,hashcash}.
Нода решает заголовок блока над свежей вершиной на облегчённую цель и прикладывает его как PoW-билет допуска.
В коинбейсе билета закреплены личность, проект, роль, эпоха и дайджест политики, поэтому билет относится к одной заявке.
Если заголовок попал и в цель блока, это блок.
Билет служит шаром: работа допуска охраняет цепочку, а свежая вершина служит вызовом без внешнего маяка случайности, и решить билет до её появления нельзя.

\section{Финальность: штраф поздней ветки}
\label{sec:finality}

Молодая CPU-сеть хэшрейтом не защитится: большинство можно арендовать.
Время арендовать нельзя, и на этом строится выбор ветки.
Правило продолжает штраф за задержанную публикацию, предложенный для Horizen \cite{horizen}, и MESS, который Ethereum Classic ввёл после атак 2020 года \cite{mess} и позднее отключил \cite{messoff}.

Узел помнит монотонное время первого получения каждого блока.
Пусть конкурирующая ветка расходится с цепочкой узла в точке, после которой первый блок он получил $\Delta$ назад.
Узел переходит на неё, только если её работа после точки расхождения больше работы его текущей цепочки после той же точки в $P(\Delta)$ раз:
\[
P(\Delta)=
\begin{cases}
1, & \Delta\le 10~\text{мин},\\
\min\!\bigl(32,\ 2^{(\Delta-10~\text{мин})/1~\text{ч}}\bigr), & 10~\text{мин}<\Delta<24~\text{ч},\\
\infty, & \Delta\ge 24~\text{ч}.
\end{cases}
\]
Первые 10 минут штрафа нет, поэтому обычные форки он не затрагивает.
Дальше штраф удваивается каждый час до потолка 32, а через 24 часа реорганизаций нет.
Параметры служат кандидатами, и симуляция проверяет их вместе с ASERT и правилами времени.

Атакующий с долей хэшрейта $\alpha$ производит за то же время в $k=\alpha/(1-\alpha)$ раз больше работы, чем остальная сеть.
Пусть получатель ждал $w$ после первого блока с платежом и только тогда выдал товар.
Тайная ветка, показанная после этого, в среднем побеждает только при $k>P(w)$, то есть при $\alpha>P(w)/(1+P(w))$.
Без штрафа атакующий с долей больше половины рано или поздно побеждает при любом ожидании \cite{bitcoin,rosenfeld}.
Со штрафом придерживать блоки дольше форы бессмысленно.
Численные значения приведены в таблице~\ref{tab:penalty}.

\textbf{Слабая субъективность.}
Плата за отказ от поздних реорганизаций называется слабой субъективностью \cite{weaksubj}.
Новый узел стартует с подписанного checkpoint и снимка состояния (UTXO, имён и личностей) и сверяет снимок с корнем в заголовке; сырые социальные сегменты для этого не нужны.
Checkpoint подписывают хэш-ключи сопровождающих, и он публикуется в релизе и на зеркалах.
Узел, вернувшийся после долгого отсутствия, сверяет вершины последних суток у нескольких пиров.
Разделение сети дольше суток разрешается вручную по опубликованной процедуре.

\textbf{Без второго консенсуса.}
Финальность на кворумах по образцу ChainLocks \cite{chainlocks} исключена: второй консенсус на пороговых BLS противоречил бы хэш-профилю, а его задачу решает штраф поздней ветки.
Три усиления ядро не включает, и дизайн держит их в резерве.
Последовательная композиция двух CPU-функций с разными узкими местами ограничила бы выигрыш от прорыва в одной из них примерно двукратным.
Нотариальная отметка вершины в Bitcoin помогла бы новым узлам и разбору разделений, а выбор ветки от неё бы не зависел.
STARK-агрегация подписей \cite{stark}, тоже на одних хэшах, сняла бы цену их размера.

\section{Лёгкие клиенты и индексаторы}
\label{sec:light}

\textbf{Заголовки от checkpoint.}
Заголовок с 64-байтными обязательствами весит около 350 байт, а проверка RandomX на телефоне тяжела.
Поэтому лёгкий клиент проверяет заголовки от последнего подписанного checkpoint и всю историю не хранит.
Его доверие к истории до checkpoint сводится к подписантам checkpoint, а после него к работе в заголовках.
На iOS клиент работает в тонком режиме, без архивирования.

\textbf{Живая ячейка по корню.}
Корень состояния личностей и имён в каждом заголовке позволяет лёгкому клиенту проверить по доказательству против этого корня, что ячейка выдавшего была жива на высоте записи.
Вместе с путём Меркла пачки и подписью это закрывает три первых шага проверки аттестации (раздел~\ref{sec:attest}) на устройстве с одними заголовками.

\textbf{Фильтры.}
Лёгкий клиент читает компактные фильтры блоков по образцу BIP 157/158 \cite{bip157,bip158}; обязательство фильтра лежит в заголовке.
Фильтр для социальных записей строится по четырём ключам: тип записи, автор, ссылка на запись и тема.
Так клиент ищет записи выбранного автора или ответы на известную запись, не отдавая серверу список подписок.
Фильтр указывает на возможное совпадение, а полученные данные проверяются по заголовку.
Запись, которую больше никто не раздаёт, фильтр не вернёт: потеря сегментов остаётся потерей истории.

\textbf{Проверяемый индекс.}
Индексатор отдаёт производное состояние через GraphQL с подписками: ленты по личности и по имени, граф подписок, деревья ответов, разрешение указателей, репутацию по опубликованной формуле и доказательства включения.
На заданной высоте он публикует хэш своего состояния.
Клиент сравнивает ответы нескольких индексаторов на одной высоте при одинаковых правилах вывода.
Совпадение подтверждает согласованность канонического вывода, а правильности произвольной рекомендации не доказывает.
Точную схему задаёт спецификация индексатора.

\section{Как рантайм и протокол используют цепочку}
\label{sec:use}

Рантайм строится в расчёте на Swarma Chain, и через цепочку связаны личность агента, его свидетельства и переносимая память (табл.~\ref{tab:use}).

\begin{table}[t]
\centering\small
\caption{Что слои Swarma берут из цепочки и где проходит граница.}
\label{tab:use}
\footnotesize
\begin{tabular}{L{2.5cm}L{7.9cm}L{5.2cm}}
\toprule
Слой & Что берёт из цепочки & Граница \\
\midrule
Swarma Runtime & Ячейки агентов с ключами у ключника; журнал ключевых событий в реестре владельца с тем же идентификатором; ссылки памяти на ячейки & Происхождение ключа не доказывает честности агента \\
Swarma Protocol & Личности нод и подписантов манифеста; пачки аттестаций исхода; хэш блока для выбора ревьювера; PoW-билеты допуска; закрепление голов журнала & Контроль над внешней сущностью не доказывает независимости ошибок \\
Swarma Chain & Ячейки, имена и корень их состояния в каждом заголовке & Общие записи не заменяют политику допуска принимающего проекта \\
\bottomrule
\end{tabular}
\end{table}

\textbf{Рантайм.}
Ключник задаёт личность агента стабильным путём вида назначение, продукт, проект, агент; путь отличается от домашнего каталога и от UID пользователя ОС.
Учётку можно пересоздать, а агента перенести на другую машину, и личность при этом не меняется.
Ключи ячейки ключник выводит хэш-функцией от пути и номера ротации и подписывает по запросу агента.
Инъекция в промпт действует возможностями процесса агента, поэтому секрет не лежит в доступном ему файле.
Честности запросов подписи ключник при этом не обеспечивает: сообщение, подписанное захваченным агентом, остаётся неправильным сообщением.
Реестр владельца ведёт журнал ключевых событий до публикации в цепочку (раздел~\ref{sec:identity}).
Память агентов называет издателя, субъекта и отправителя записей идентификаторами ячеек, а представительство опирается на грант в реестре владельца; недоступный реестр никогда не считается выданным грантом.
Поэтому переезд агента на другую машину или в другой рантайм не рвёт ни истории, ни памяти: меняются ключи и место исполнения, а ячейка остаётся.

\textbf{Протокол.}
Личность ноды задаёт её ячейка, а короткоживущие ключи ноды удостоверяет эта личность.
Ревизии подписываются ключом ноды в формате SSH-подписи git, и ключ допустим, только если его удостоверила подтверждённая личность и нода назначена на задачу.
Аттестации исхода из журнала проекта выпускаются пачками в формате цепочки, и их корни публикуются в Swarma Chain.
Манифест федерации подписывает личность цепочки, а доказательства идентичности организации привязываются к этой личности аттестациями.
Головы журналов протокола закрепляются в цепочке, и записи на Ed25519 сохраняют доказательную силу на дату наблюдения.

Выбор ревьювера опирается на хэши цепочки.
Кор-нодой называется узел, который держит журнал проекта на кворуме.
Протокол замораживает множество допустимых ревьюверов до появления хэша кандидата.
Случайность равна хэшу раскрытий луковиц хэшей кор-нод вместе с хэшем блока цепочки, вышедшего после фиксации кандидата.
Ревьювер вычисляется из идентификатора задачи, хэша кандидата и случайности, и валидатор пересчитывает результат.
Публичная точка присоединения принимает заявку с PoW-билетом (раздел~\ref{sec:consensus}) и ключа подписи допуска не держит: точка ставит заявку в очередь, а допуск подписывает проект.

\section{Анализ безопасности}
\label{sec:security}

Для каждой цели раздела~\ref{sec:model} мы называем аргумент, допущения и предел.

\textbf{Ц1: подделка состояния.}
Каждое изменение денег, имён и ячеек требует пары подписей SLH-DSA, а ключи для проверки берутся из самого состояния: из условия траты выхода или из поля ячейки.
Противнику П1, сломавшему кривые, пейринги или решётки, подделывать нечего, потому что ни одна проверка консенсуса на них не опирается.
Противнику, сломавшему одно семейство хэшей, нужна вторая половина пары, которая держится на другом семействе, а обязательства тем же способом требуют коллизии в обоих семействах.
Противник П2 с большинством хэшрейта ничего не подделывает, он только упорядочивает.
Допущение: Д1.
Предел: поломка обоих семейств хэшей ломает профиль целиком.

\textbf{Ц2: перехват личности.}
Вор текущих ключей (П3) подписывает записи, но потратить ячейку не может: ротация требует следующих ключей, спрятанных за хэшем, а трата родителем требует ключей ячейки владельца.
Всё, что вор подписал до ротации, проходит проверку подписи, и владелец оспаривает такие записи возражениями.
Вор мастер-секрета знает и следующие ключи агентов и может провернуть ячейку.
Тогда работает второй путь: собственная ротация не меняет ссылку на родителя, поэтому ячейка владельца тратит текущую версию ячейки агента в любой момент и задаёт ключи, которых вор не знает.
Если трата вора и трата родителя конкурируют за одну версию, блок включает только одну из них; после включения траты вора родитель тратит следующую версию.
Допущение: Д6.
Предел: устойчивое большинство может задерживать трату родителя, а потеря и ключа владельца, и ключа восстановления делает личность невосстановимой.

\textbf{Ц3: одна версия.}
Двуличный владелец (П4) в реестре одного владельца может показать разные истории разным проверяющим.
В цепочке ячейка лежит в общем состоянии, корень которого входит в каждый заголовок, поэтому два честных узла на одной ветке видят одну версию.
Разные ветки возможны до созревания: противник П2 может реорганизовать свежую историю в пределах штрафа, но не после 24 часов.
Допущения: Д3, Д5.
Предел: узел под eclipse-атакой (П7) видит ветку противника до получения честной, а узел, вернувшийся после суток отсутствия, полагается на сверку вершин у нескольких пиров.

\textbf{Ц4: проверка аттестации.}
Подпись проверяется по ключам выдавшего на высоте записи, включение по пути Меркла до корня пачки и заголовка, статус по корню состояния личностей в заголовке.
Ни индексатор (П6), ни выдавший в этих трёх шагах не участвуют.
Индексатор может скрыть запись или подставить другую формулу.
Скрытие ограничено тем, что субъект носит доказательства сам, а подмену формулы обнаруживает пересчёт, потому что вывод детерминирован.
Против фабрики личностей (П5) действует вес.
Кластер, который аттестует сам себя, получает вес только от внешнего притока, а каждая личность и каждая запись стоят платы.
Предел: доказательство включения не доказывает полноты выдачи, и отсутствие записи в ответе не значит, что её не было.
Корректная подпись не делает истинным утверждение внутри свидетельства.
При пропорциональном вознаграждении полная устойчивость к множественным личностям недостижима \cite{sybilproof}, поэтому репутация снижает риск сговора и создания личностей, но не исключает его.

\textbf{Ц5: двойная трата.}
При доле противника $\alpha$ и ожидании получателя $w$ тайная ветка побеждает только при $\alpha>P(w)/(1+P(w))$ (табл.~\ref{tab:penalty}).
Штраф считается по монотонным часам узла и от меток блоков не зависит, поэтому сдвиг меток его не обходит.
Метки влияют только на ASERT, где их ограничивают строгая монотонность и допуск 300 секунд в будущее.
Стратегии эгоистичного майнинга \cite{selfish}, которые придерживают блоки дольше 10 минут, требуют кратного перевеса.
Допущения: Д2, Д3, Д4.
Предел: в первые 10 минут $P=1$, и действует обычная гонка ветвей \cite{rosenfeld}, поэтому получатель крупного платежа ждёт часы.
Узлы с разным временем первого получения блоков могут по-разному оценить одну конкурирующую ветку у границы порога, и их взгляды расходятся до явного перевеса одной из веток; разделение дольше суток разрешается вручную.

\textbf{Ц6: цена записи.}
Минимальная базовая плата привязана к награде так, что полный по цели блок стоит не меньше $R/16$.
Заливка поэтому стоит денег с первого байта, а при переполнении блоков базовая плата растёт на 12,5\,\% за блок.
Социальный сегмент занимает не больше половины веса, и платежи, имена и события личностей конкурируют за свою половину.
Предел: подпись спама не останавливает, потому что спамер тоже подписывает свои записи, и долларовую цену записи протокол не обещает.
Видимость и хранение принятого спама решают клиенты, кураторы и архивисты.

\textbf{Ц7: дата наблюдения.}
Голова журнала протокола или реестра, закреплённая в блоке, доказывает, что журнал в этом состоянии существовал не позже этой высоты, и это доказательство держится только на хэшах.
После поломки кривых противник подделает подписи живых ключей, но не перенесёт поддельную запись в прошлое закреплённой головы.
Допущения: Д1, Д5.
Предел: записи между поломкой и отказом политики от класса ключей по подписи не проверить, а короткоживущие листья противник подделывает на время их действия.

\textbf{Случайность выбора.}
Хэш раскрытий луковиц хэшей и хэш последующего блока никто не подбирает после фиксации кандидата.
Предел: участник, раскрывающий последним, может не раскрыть прообраз, а майнер может отбросить найденный блок с неудобным хэшем ценой награды за него.
Каждое такое действие даёт противнику одну новую попытку и не позволяет назначить выбранного им ревьювера из замороженного множества.

\textbf{Остаточный риск.}
Устойчивое большинство может часами задерживать записи.
Отдельный узел можно изолировать eclipse-атакой.
Атака на небольшую сеть дешевле, чем на Bitcoin, и параметры финальности рассчитаны именно на арендованное большинство.

\section{Анализ параметров}
\label{sec:params}

Числа ниже получены расчётом из кандидатных параметров дизайна.
Окончательные значения генезиса задаёт моделирование сложности, размера блока, эмиссии, штрафа поздней ветки и цены атаки.

\textbf{Подписи и пропускная способность.}
Парная подпись весит 15\,712 байт.
Блок по цели $T_0=512$~КБ вмещает около $512\cdot1024/15\,712\approx33$ подписанных объектов, около одного объекта в две секунды, а блок наибольшего размера $2T_0$ вмещает около 66.
При полном социальном сегменте телу остаётся половина веса, около 16 объектов на блок.
Поэтому события личностей, имена и платежи редки по замыслу, а массовые записи идут пачками.

\textbf{Пачки.}
В пачке из $N$ записей на одну запись приходится $15\,712/N$ байт подписи и путь в $64\lceil\log_2 N\rceil$ байт (табл.~\ref{tab:batch}).
Начиная примерно с $N=256$, путь весит больше доли подписи, и дальнейший рост пачки почти не снижает цену записи.

\begin{table}[t]
\centering\small
\caption{Доля подписи и длина пути на одну запись пачки, байт.}
\label{tab:batch}
\footnotesize
\begin{tabular}{lrrrrr}
\toprule
Записей в пачке $N$ & 16 & 64 & 256 & 1024 & 4096 \\
\midrule
Доля парной подписи, $15\,712/N$ & 982 & 246 & 61 & 15 & 4 \\
Путь Меркла, $64\lceil\log_2N\rceil$ & 256 & 384 & 512 & 640 & 768 \\
\bottomrule
\end{tabular}
\end{table}

\textbf{Цена записи.}
Минимальная плата за байт равна $f_{\min}=R/(16\,T)$.
Запись в 1~КБ при цели 512~КБ стоит не меньше $R/8192$, то есть около 0,0078 монеты при начальной награде в 64 условные монеты.
Событие личности с парной подписью весит больше 15~КБ и стоит не меньше примерно $R/530$, около 0,12 монеты при той же награде.
Это нижняя граница цены каждой новой личности, и при переполненных блоках она растёт вместе с базовой платой.

\textbf{Рост цели.}
Цель растёт не быстрее чем вдвое в год, а потолок равен $4~\text{МБ}\cdot2^{y/4}$.
При непрерывном наибольшем росте от 0,5~МБ цель $0{,}5\cdot2^{y}$~МБ догоняет потолок при $2^{3y/4}=8$, то есть к концу четвёртого года, на 8~МБ.
Дальше её ограничивает потолок, который растёт вдвое за четыре года, примерно на 19\,\% в год, а наибольший блок в этот момент равен 16~МБ.

\textbf{Цепочка заголовков.}
При 60-секундном блоке за год выходит $525\,600$ блоков, и заголовки по 350 байт занимают около 184~МБ в год.
Лёгкий клиент хранит только заголовки от последнего checkpoint.

\textbf{Двойная трата.}
Таблица~\ref{tab:penalty} даёт штраф и наименьшую долю хэшрейта, при которой тайная ветка в среднем побеждает.
Без штрафа порог равен половине при любом ожидании.

\begin{table}[t]
\centering\small
\caption{Штраф поздней ветки и наименьшая доля хэшрейта атакующего, $\alpha>P/(1+P)$.}
\label{tab:penalty}
\footnotesize
\begin{tabular}{lrrrrrr}
\toprule
Ожидание получателя $w$ & 10 мин & 1 ч & 2 ч & 3 ч & 5 ч 10 мин & 24 ч \\
\midrule
$P(w)$ & 1 & 1,78 & 3,56 & 7,13 & 32 & $\infty$ \\
Нужная доля $\alpha$ & 50,0\,\% & 64,1\,\% & 78,1\,\% & 87,7\,\% & 97,0\,\% & невозможно \\
\bottomrule
\end{tabular}
\end{table}

\textbf{Сложность.}
Для кандидата $\tau=2$~ч предварительные оценки дают шум сложности около 5\,\%, приспособление к десятикратному росту хэшрейта примерно за четыре часа и медленные блоки около шести с половиной часов после ухода такого майнера.
Порядок величины виден из формулы ASERT: чтобы цель изменилась в десять раз, отставание от расписания или опережение должно достичь $\tau\log_2 10\approx6{,}6$~ч.
Симуляция проверяет скачки хэшрейта, уход крупного майнера и взаимодействие пересчёта с метками времени.

\textbf{Эмиссия.}
Таблица~\ref{tab:emission} сравнивает три графика для минутного блока; монеты условные.
Для выбранного плавного спада награда равна $R(y)=64\cdot2^{-y/3}$ до пятнадцатого года, а выпуск до хвоста равен
\[
\int_0^{15} 64\,B\,2^{-y/3}\,dy = 64\,B\,\frac{3}{\ln 2}\,\bigl(1-2^{-5}\bigr)\approx 141~\text{млн},
\qquad B=525\,600~\text{блоков в год}.
\]
Половина этого выпуска приходится на первые три года ($51{,}6\,\%$), больше 90\,\% на первые десять.
Хвост в 2 монеты за блок даёт $1\,051\,200$ монет в год, около 0,75\,\% накопленного к его началу предложения, и дальше эта доля уменьшается.
График с халвингами выпускает до хвоста $50\cdot4B\cdot(1+\frac12+\frac14+\frac18)\approx197$ млн монет.
Сжигаемая базовая плата и минимальная плата, привязанная к награде, связывают график эмиссии с правилом размера блока, поэтому их числа моделируются вместе.

\begin{table}[t]
\centering\small
\caption{Графики эмиссии для минутного блока, условные монеты.}
\label{tab:emission}
\footnotesize
\begin{tabular}{L{2.9cm}L{4.0cm}L{3.0cm}r L{3.5cm}}
\toprule
График & Начало и спад & Хвост & До хвоста & Компромисс \\
\midrule
Халвинги и хвост & 50 за блок, вдвое каждые 4 года & 3,125 с 16-го года & 197 млн & Резкое падение дохода майнеров \\
Плавный спад и хвост (выбран) & 64 за блок, плавно вдвое каждые 3 года & 2 с 15-го года & 141 млн & Нужна целочисленная экспонента \\
Быстрый спад и хвост & 200 за блок, помесячное снижение & 2 примерно с 7-го года & 150 млн & Преимущество ранних участников \\
\bottomrule
\end{tabular}
\end{table}

\section{Смежные работы}
\label{sec:related}

\textbf{Цепочки с доказательством работы.}
Bitcoin \cite{bitcoin} ввёл цепочку наибольшей работы и UTXO, Ethereum \cite{ethereum,ethyellow} ввёл общее состояние с контрактами, а анализ протокола Bitcoin \cite{backbone} задаёт свойства общего префикса и качества цепочки при честном большинстве.
Swarma Chain берёт UTXO и наибольшую работу и меняет правило выбора ветки: поздняя ветка платит за задержку кратным перевесом работы.
Эгоистичный майнинг \cite{selfish}, eclipse-атаки \cite{eclipse} и анализ двойной траты \cite{rosenfeld} задают атаки, против которых рассчитаны параметры.
Штраф поздней ветки продолжает предложение Horizen \cite{horizen} и MESS \cite{mess,messoff}, а checkpoint следует рассуждению о слабой субъективности \cite{weaksubj}.
ChainLocks \cite{chainlocks} дают финальность кворумом мастернод на пороговых BLS, и хэш-профиль её исключает.
Объединённый майнинг разобран в \cite{mergedmining}, а RandomX \cite{randomx}, P2Pool \cite{p2pool}, ASERT \cite{asert} и хвостовая эмиссия Monero \cite{tail} служат образцами компонентов.

\textbf{Хэш-подписи.}
Одноразовые подписи Лэмпорта \cite{lamport79} и деревья Меркла \cite{merkle89} лежат в основе хэш-подписей.
XMSS и LMS \cite{xmss,rfc8391,rfc8554} компактнее, но держат состояние ключа, и NIST задаёт для них отдельные требования к управлению состоянием \cite{sp800208}; QRL \cite{qrl} строит цепочку на XMSS.
SPHINCS, SPHINCS$^+$ и стандарт SLH-DSA \cite{sphincs,sphincsplus,fips205} убирают состояние ценой размера.
Swarma Chain берёт SLH-DSA парой из двух семейств хэшей, чтобы поломка одного семейства не давала подделки.
STARK \cite{stark} дают агрегацию на тех же предположениях.
Угрозу кривым задаёт алгоритм Шора \cite{shor}, а алгоритм Гровера \cite{grover} даёт перебору прообразов хэш-функции лишь квадратичное ускорение.

\textbf{Личности и прозрачность.}
KERI \cite{keri} вводит самосертифицирующие идентификаторы, журнал ключевых событий и предротацию.
Swarma переносит журнал в состояние консенсуса, где раздвоение исключает сама цепочка, и добавляет внешнего родителя как второй путь траты.
Certificate Transparency \cite{ct,ct2}, CONIKS \cite{coniks} и Sigstore \cite{sigstore} делают раздвоение журналов обнаружимым для аудиторов.
Иерархическая деривация \cite{bip32,slip10} задаёт ключи каналов рантайма.

\textbf{Имена.}
Namecoin \cite{namecoin} поставил имена в цепочку, а его эмпирическое изучение \cite{namecoinstudy} показало, что пространство имён без сдерживания сквоттинга захватывается.
ENS \cite{ens} задаёт имена контрактом поверх Ethereum.

\textbf{Репутация, цена работы и лёгкие клиенты.}
EigenTrust \cite{eigentrust} вычисляет глобальное доверие как неподвижную точку графа, а атака множественными личностями \cite{sybil} и пределы устойчивых к ней механизмов \cite{sybilproof} задают границы репутации.
Цена работы против спама восходит к \cite{dworknaor,hashcash}, и PoW-билеты Swarma делают эту работу шарами майнинга.
Рынок платы следует EIP-1559 \cite{eip1559}, фильтры лёгкого клиента следуют BIP 157/158 \cite{bip157,bip158}.
FlyClient \cite{flyclient} сокращает проверку заголовков вероятностной выборкой; Swarma Chain вместо этого начинает проверку с подписанного checkpoint.

\section{Заключение}

Swarma Chain даёт рою агентов общее состояние там, где кончается знание одного владельца.
Личность агента переживает смену ключей, машин и рантаймов, аттестация проверяется по корням заголовка без доверия индексатору, а имя и расчёты живут в том же состоянии.
Цепочка заверяет записи, как нотариус: самоаутентичность держится на одних хэш-функциях, а доказательство работы и штраф поздней ветки отвечают только за порядок.
Цена этого выбора названа прямо: подпись весит больше 15~КБ, блок вмещает десятки подписанных объектов, расчёты медленные, а новый узел доверяет checkpoint.

Открытыми точками дизайна остаются числа генезиса (цель и темпы размера блока, минимальная плата, полупериод ASERT, параметры штрафа и графика эмиссии), которые задаёт моделирование, константы формулы репутации, параметры оплаты сервисных узлов и связь денег с репутацией.
Каждая из них задаёт параметр; граница консенсуса, профиль хэшей, ячейка личности и правило выбора ветки от них не зависят.
Ячейка доказывает происхождение ключа, аттестация доказывает исход процедуры, а решение о допуске остаётся у проекта.

{\footnotesize
\bibliographystyle{abbrv}
\setlength{\itemsep}{0pt}
\bibliography{references}
}

\end{document}
